\section{Conclusiones}

El objetivo principal de este trabajo era desarrollar una plataforma funcional y validada de
comunicación con periféricos serie sobre el MPSoC Zynq UltraScale+ ZCU102, que sirviera de base
para el ordenador de a bordo del proyecto LINCE. Ese objetivo se desglosó en la
sección~\ref{objetivos} en seis objetivos específicos, y las conclusiones se ordenan aquí
siguiéndolos uno a uno, para poder declarar en cada caso qué se ha conseguido y con qué evidencia.

\textbf{El primer objetivo, diseñar un transceptor serie configurable en VHDL, se ha cumplido.}
El bloque desarrollado opera sobre los estándares RS422 y RS485 y admite la configuración completa
del formato de trama en tiempo de ejecución: velocidad entre 50 y 4\,000\,000 de baudios sobre una
tabla de 54 entradas, de 5 a 9 bits de datos, cinco modos de paridad, tres campos de parada
(incluido el de 1,5 bits) y los dos órdenes de bit. Su base de tiempos, un NCO de 32 bits validado
velocidad a velocidad con un \textit{testbench} dedicado, presenta un error nulo o inferior a
1~ppm en la mayoría de las velocidades estándar y de 120~ppm en el peor punto del rango, dos
órdenes de magnitud por debajo de la tolerancia habitual de un receptor UART. El canal completo
supera además el barrido exhaustivo de las 150 configuraciones posibles, 450 tramas comparadas dato
a dato sin ningún error (sección~\ref{sec:verificacion_canal}), y catorce instancias del bloque
funcionan de forma simultánea sobre la placa.

\textbf{El segundo objetivo, diseñar la arquitectura de transporte entre la lógica programable y
el sistema de procesamiento, es el de mayor recorrido del trabajo y se ha cumplido con la
arquitectura MCDMA que gobierna los catorce canales del sistema final.} No se seleccionó una
alternativa sobre el papel: se implementaron y midieron
\textbf{tres arquitecturas completas} (AXI GPIO, catorce AXI DMA y un AXI MCDMA con puente VHDL
propio) sobre el mismo transceptor y con el mismo programa de pruebas. La de partida, basada en AXI
GPIO, interrumpe al procesador una vez por byte, se queda en el 30\,\% del techo físico de la línea
y apenas mejora al subir la velocidad, porque el cuello de botella está en la CPU. La solución final
baja a \textbf{3 interrupciones por kilobyte}, alcanza el \textbf{99\,\% del techo físico} y escala
linealmente hasta los 345~kB/s a 4~Mbaudios. La variante de catorce DMA funcionó en placa, pero se
descartó por recursos: 2,8 veces más lógica, veintiocho líneas de interrupción donde el MPSoC solo
ofrece ocho y un área incompatible con una futura triplicación TMR. Ese límite del silicio justificó
el diseño del puente propio.

\textbf{El tercer objetivo, un driver completo en C sobre RTEMS 7, se ha cumplido en sus tres
versiones.} El driver abstrae por completo el acceso al hardware tras una API pública de seis
funciones, descubre en arranque cuántos transceptores hay instanciados en el \textit{bitstream}
(de modo que el mismo binario sirve para cualquier configuración, de un canal a catorce) y gestiona
la comunicación orientada a interrupciones sin bloquear al procesador. La separación entre la ISR
maestra y las tareas \textit{worker} mantiene mínimo el tiempo en contexto de interrupción y
preserva el determinismo propio de un sistema de tiempo real. Se escribió una versión del driver
para cada arquitectura de transporte, lo que permitió que la comparativa del segundo objetivo
midiera transportes y no implementaciones desiguales.

\textbf{El cuarto objetivo, diseñar y fabricar una placa que ejercitara los catorce transceptores
sobre topologías realistas, se ha cumplido y la placa se ha caracterizado eléctricamente.} Sus
catorce canales se reparten en un bus RS485 multipunto de siete nodos y dos buses RS422
\textit{master/slave}, configurables por \textit{jumper} sin recablear. La batería de nueve tests
recorre la polarización de reposo, el mapa de conectividad real, la integridad del bus multipunto,
la velocidad máxima de cada bus, la diafonía, la colisión y la estabilidad sostenida, y la placa
los supera todos: cero bytes espurios en reposo, cero fugas entre buses, recuperación limpia tras
una colisión y tasa de error nula en los dos buses RS422 tras diez segundos de tráfico continuo.
La campaña reveló además que el \textbf{limitador de \textit{slew rate}} de los transceptores, activo
en la configuración por defecto del software, fijaba el techo de los buses entre 460~kbaudios y
1~Mbaudios. Con él desactivado, los tres buses alcanzan \textbf{4~Mbaudios} con tasa de error nula,
sin modificar el cobre ni resintetizar la lógica.

\textbf{El quinto objetivo, diseñar y fabricar las dos placas de expansión bajo especificación de
Indra, se ha cumplido.} La placa CDHS expone CAN redundante, tres canales RS422/RS485, cuatro
líneas de PWM para calentadores y un ADC de termistores sobre seis conectores D-Sub-9; la placa
AOCS, los canales serie, el control de motores por PWM y dos enlaces SpaceWire sobre ocho. Las tres
tarjetas del trabajo se fabricaron en el propio laboratorio mediante soldadura por reflujo con
\textit{stencil}, y las tres se montaron sin retrabajos por haberse diseñado directamente contra
las capacidades de proceso del fabricante. La restricción de 1,8~V del carril VADJ del FMC
condicionó la selección de componentes y obligó a descartar las piezas con herencia de vuelo, una
decisión que se declara explícitamente: las tarjetas son hardware de validación funcional, no
hardware de vuelo. La placa AOCS presenta además un error de conexión en el eje Z del bloque de
motores, corregido en la revisión V2 del esquemático.

\textbf{El sexto objetivo, integrar y validar el sistema completo, se ha cumplido salvo en una
interfaz.} La validación experimental confirma el funcionamiento de las interfaces RS422 y RS485 de
las placas CDHS y AOCS, del bus CAN de la CDHS (26 de 26 pruebas superadas, incluida la
comprobación del efecto de las resistencias de terminación sobre la integridad de la señal
diferencial), del ADC por SPI, cuyas lecturas siguen linealmente la tensión aplicada, y de las
señales de PWM de ambas placas. La única interfaz que quedó sin validar es \textbf{SpaceWire}, por
una razón declarada desde el inicio: la placa expone la capa física LVDS, pero ejercitarla exige
implementar en la lógica programable la codificación \textit{Data-Strobe}, la máquina de estados
del enlace y los niveles superiores del estándar, un desarrollo que por sí solo excede el alcance
de este trabajo.

En términos de volumen, el trabajo ha producido \textbf{20\,188 líneas de código} de autoría propia,
sin contar el código generado ni el de terceros, desglosadas en la tabla~\ref{tab:loc} del
Anexo~\ref{anx:rtl}. El código de verificación pesa prácticamente lo mismo que el hardware que
verifica, y la automatización del flujo, con un 40\,\% del total, es la mayor partida individual.

En conjunto, los seis objetivos específicos quedan cubiertos y el objetivo principal, con ellos: el
resultado es una plataforma validada sobre medidas propias, con una arquitectura de transporte
escogida a partir de ellas y no sobre estimaciones de catálogo. Esa arquitectura reduce la carga de
interrupciones a 3 por kilobyte y ocupa menos lógica que la alternativa de catorce motores DMA, lo
que deja margen de CPU y de área para el software de vuelo. El trabajo constituye así una
contribución directa al proyecto LINCE y deja al laboratorio B105 una base sobre la que construir
el subsistema de comunicaciones del ordenador de a bordo del satélite.
